\documentclass[a4paper, 12pt]{extarticle}


%%%%%%%%%%%%%%%%%%%%%%%%%%%%%%
% PACKAGES
%%%%%%%%%%%%%%%%%%%%%%%%%%%%%%

%!TEX encoding = UTF8

%%%%%%%%%%%%%%%%%%%%%%%%%%%%%%%%%
% PACKAGE IMPORTS
%%%%%%%%%%%%%%%%%%%%%%%%%%%%%%%%%


\usepackage[french]{babel}
\usepackage[T1]{fontenc}
\usepackage[utf8]{inputenc}

\usepackage[
	margin=2cm,
	right=2.2cm,
	footskip=1.5cm,
]{geometry}

%ams
\usepackage{amsmath,amsfonts,amsthm,amssymb,mathtools}

\usepackage{bookmark}
\usepackage{enumitem}
\usepackage[most,many,breakable]{tcolorbox}
\usepackage{varwidth,etoolbox}
\usepackage[makeroom]{cancel}
\usepackage{xcolor}
\usepackage{multicol,array}
\usepackage[ruled,vlined,linesnumbered]{algorithm2e}

\usepackage{caption, subcaption}

\usepackage{marginnote}

\usepackage{makecell} % pour \thead 

%index
\usepackage[texindy]{imakeidx}
\makeindex[
	columns=2, 
	title=Index, 
	options= -L french,
	%intoc,
]
\indexsetup{level=\chapter}

%virgules
\usepackage{icomma}

% for striked out implies sign (\centernot\implies)
\usepackage{centernot}

% roman numerals for \section
\renewcommand{\thesection}{\Roman{section}} 

% tikz
\usepackage{tikz}

% package systeme pour les systèmes d'équations bien alignés
\usepackage{systeme}
%\sysalign{r,r}
%\syseqspace{3pt}
%\syssignspace{3pt}

\newcommand{\sys}[2]{\systeme{#1{,}, #2.}}

% date
\usepackage{advdate}

\pagestyle{fancy}
\fancyhead[L]{Université de Mayotte}
\fancyhead[C]{\bf Série 1}
\fancyhead[R]{\today}

%%%%%%%%%%%%%%%%%%%%%%%%%%%%%%
% SELF MADE COLORS
%%%%%%%%%%%%%%%%%%%%%%%%%%%%%%

\input{../../../colors}

%%%%%%%%%%%%%%%%%%%%%%%%%%%%
% LETTERFONTS
%%%%%%%%%%%%%%%%%%%%%%%%%%%%

\input{../../../letterfonts}

%%%%%%%%%%%%%%%%%%%%%%%%%%%%
% MACROS
%%%%%%%%%%%%%%%%%%%%%%%%%%%%

%!TEX encoding = UTF8


%%%%%%%%%%%%%%%%%%%%%%%%%%%%
% THEOREMS
%%%%%%%%%%%%%%%%%%%%%%%%%%%%

\newcommand{\thm}[3]{\begin{theorem}[#1]\label{#3}#2\end{theorem}}
\newcommand{\cor}[3]{\begin{corollaire}[#1]\label{#3}#2\end{corollaire}}
\newcommand{\lem}[3]{\begin{lemme}[#1]\label{#3}#2\end{lemme}}
\newcommand{\prop}[3]{\begin{proposition}[#1]\label{#3}#2\end{proposition}}
\newcommand{\ex}[3]{\begin{exemple}[#1]\label{#3}#2\end{exemple}}
\newcommand{\dfn}[3]{\begin{definition}[#1]\label{#3}#2\end{definition}}
\newcommand{\qs}[2]{\begin{question}[#1]#2\end{question}}

\renewcommand\qedsymbol{$\blacksquare$}
\newcommand{\pf}[2]{\begin{preuve}[#1]#2\hfill\qedsymbol\end{preuve}}

\newcommand{\nt}[1]{\begin{remarque}#1\end{remarque}}
\newcommand{\str}[1]{\begin{strategie}#1\end{strategie}}
\newcommand{\mth}[1]{\begin{methode}#1\end{methode}}
\newcommand{\ax}[3]{\begin{axiome}[#1]\label{#3}#2\end{axiome}}
\newcommand{\notations}[1]{\begin{notation}#1 \end{notation}}
\newcommand{\nomen}[1]{\begin{nomenclature}#1 \end{nomenclature}}
\newcommand{\heur}[1]{\begin{heuristique}#1\end{heuristique}}
\newcommand{\motiv}[1]{\begin{motivation}#1\end{motivation}}


%%%%%%%%%%%%%%%%%%%%%%%%%%%%
% EXERCISES
%%%%%%%%%%%%%%%%%%%%%%%%%%%%

\newcommand{\exe}[4]{
	\begin{Exercise}[#1, label=#3]
		%\marginpar{\mbox{\scriptsize(solution p.\pageref{\ExerciseLabel-Answer})}}
		#2
	\end{Exercise}
	\begin{Answer}[ref=#3]
		#4
	\end{Answer}
}

% version 1
\newcommand{\exemulticols}[5]{
	\begin{multicols}{2}
	\begin{Exercise}[#1, label=#4]
		#2
	\end{Exercise}
	\columnbreak
		#3
	\end{multicols}
	\begin{Answer}[ref=#4]
		#5
	\end{Answer}
}

% version 2
%\newcommand{\exemulticols}[5]{
%	\begin{Exercise}[#1, label=#4]
%	\begin{multicols}{2}
%		#2
%	\columnbreak
%		#3
%	\end{multicols}
%	\end{Exercise}
%	\begin{Answer}[ref=#4]
%		#5
%	\end{Answer}
%}

\newcommand{\exeTF}[1]{}

%%%%%%%%%%%%%%%%%%%%%%%%%%%%
% OPERATORS
%%%%%%%%%%%%%%%%%%%%%%%%%%%%

% deliminators
\DeclarePairedDelimiter{\abs}{\lvert}{\rvert}
\DeclarePairedDelimiter{\norm}{\lVert}{\rVert}
\DeclarePairedDelimiter{\scal}{\langle}{\rangle}

\DeclarePairedDelimiter{\ceil}{\lceil}{\rceil}
\DeclarePairedDelimiter{\floor}{\lfloor}{\rfloor}
\DeclarePairedDelimiter{\round}{\lfloor}{\rceil}

%  - others
\DeclareMathOperator{\Lap}{\mathcal{L}}
\DeclareMathOperator{\Var}{Var} % variance
\DeclareMathOperator{\Cov}{Cov} % covariance
\DeclareMathOperator{\proj}{proj} % projection
\DeclareMathOperator{\sym}{sym} % symétrique

\DeclareMathOperator{\RE}{Re} % partie réelle
\DeclareMathOperator{\IM}{Im} % partie imaginaire

\DeclareMathOperator{\pgcd}{pgcd}
\DeclareMathOperator{\ppcm}{ppcm}

\DeclareMathOperator{\ev}{ev}

\DeclareMathOperator{\tr}{tr}
\DeclareMathOperator{\sgn}{sgn}

\DeclareMathOperator{\Frac}{Frac}

% Lois
\DeclareMathOperator{\Bern}{\text{Bern}}
\DeclareMathOperator{\Binom}{\text{Binom}}

% I prefer the slanted \leq
\let\oldleq\leq % save them in case they're every wanted
\let\oldgeq\geq
\renewcommand{\leq}{\leqslant}
\renewcommand{\geq}{\geqslant}


%%%%%%%%%%%%%%%%%%%%%%%%%%%%
% commands
%%%%%%%%%%%%%%%%%%%%%%%%%%%%

% tel que
\newcommand{\tqs}{\text{ tels que }}
\newcommand{\tq}{\text{ tq. }}
\newcommand{\et}{\text{ et }}
\newcommand{\ou}{\text{ ou }}
\newcommand{\pourtout}{\text{ pour tout }}
\newcommand{\sct}{\text{ sachant }}
\newcommand{\id}{\text{id}}
\newcommand{\apriori}{\emph{a priori}}


% ensemble avec bigl et bigr
\newcommand{\bigset}[1]{\bigl\{ #1 \bigr\}}
\newcommand{\Bigset}[1]{\Bigl\{ #1 \Bigr\}}
\newcommand{\bigpar}[1]{\bigl( #1 \bigr)}
\newcommand{\Bigpar}[1]{\Bigl( #1 \Bigr)}

% PLUS INFTY AND MINUS INFTY WITH NO SPACE
\newcommand{\pinfty}{{+}\infty}
\newcommand{\minfty}{{-}\infty}

% vecteur flèche
\renewcommand{\vec}[1]{\overrightarrow{#1}}

% vecteur pmatrix
\newcommand{\pvec}[2]{\begin{pmatrix} #1 \\ #2 \end{pmatrix}}

% vecteur norme
%\newcommand{\norm}[1]{\left\Vert #1 \right\Vert}

% point plan
\newcommand{\point}[3]{
	#1\left(#2 ; #3 \right)
}

% \smash avant \underline pour coller la ligne au mot
% attention : ça empêche les line breaks donc c'est pas super...
\let\oldunderline\underline
\renewcommand{\underline}[1]{\oldunderline{\smash{#1}}}

% emph + index
% i add underline just for now
\newcommand{\emphindex}[1]{\emph{\uline{#1}}\index{#1}}

% tableau croisé
\newcommand{\tableaucroise}[4]{
\begin{tabular}{|c|c|c|c|}
	\cline{2-4}
	\multicolumn{1}{c|}{} & #1 \\ \hline
	#2 \\ \hline
	#3  \\ \hline
	#4  \\ \hline
\end{tabular}
}

% emdash in mathmode
\newcommand{\dash}{\text{---}}

%%%%%%%%%%%%%%%%%%%%%%%%%%%%
% ENVIRONMENTS
%%%%%%%%%%%%%%%%%%%%%%%%%%%%

% two columns toc
\newcommand{\twocolstableofcontents}{
	\renewcommand*\contentsname{}
	\begin{center}  
		{\Large \bfseries TABLE DES MATIÈRES \bigskip}    
	\end{center}
	\hrule
	\begin{multicols}{2}\setlength{\columnseprule}{1pt}
		\titleformat{\chapter}{}{}{0pt}{} 
		\titlespacing*{\chapter}{0pt}
		{*7} %  % vertical space before the title <<<<<<<<<<
		{*-15}  %  idem after title (in ex units + glue) <<<<<<<<<<<
		\tableofcontents
	\end{multicols}
	\hrule
}

% python minted
\newcommand{\python}[1]{
\inputminted[
		linenos,
		gobble=0,
		breaklines=true, % otherwise it breaks for no apparent reason?
		breakafter=,,
		fontsize=\normalsize,
		numbersep=8pt,
		tabsize=4, % tab ident = 4 spaces
		fontfamily=courier, %important pour les signes <, >
]{python}{python/#1.py}
}



%%%%%%%%%%%%%%%%%%%%%%%%%%%%
% PACKAGE-SPECIFIC ENVIRONMENTS
%%%%%%%%%%%%%%%%%%%%%%%%%%%%

% global settings for the first level in enumerate environments
\setlist[enumerate, 1]{%
	align = left, labelwidth = \parindent%, labelsep = 0pt, leftmargin = \parindent
}


%%%%%%%%%%%%%%%%%%%%%%%%%%%%%%
% EXERCISES
%%%%%%%%%%%%%%%%%%%%%%%%%%%%%%

\usepackage{ifthen}
\renewcommand{\ExerciseHeader}{
	\tikz[baseline=(R.base)]\node[draw,rectangle, thick, inner sep=2pt](R) {\textbf{\hyperref[\ExerciseLabel-Answer]{\theExercise}.}};\!
	\ifnum\ExerciseDifficulty=0
	\else
		(\theExerciseDifficulty)
	\fi
}
\renewcommand{\DifficultyMarker}{$\star$}
\renewcommand{\AnswerHeader}{
	\tikz[baseline=(R.base)]\node[draw,rectangle, thick, inner sep=2pt](R) {\textbf{\theExercise.}};\!
}
\renewcommand{\exe}[4]{
	\begin{Exercise}[#1, label=#3]
		\if\relax\detokenize\expandafter{\ExerciseTitle}\relax
		%\marginpar{[Bonus]}
		\else
		\marginpar{\mbox{[\ExerciseTitle]}}
		\fi
		#2
	\end{Exercise}
	\begin{Answer}[ref=#3]
		#4
	\end{Answer}
}
\renewcommand{\exemulticols}[5]{
	\begin{multicols}{2}
	\begin{Exercise}[#1, label=#4]
		\if\relax\detokenize\expandafter{\ExerciseTitle}\relax
		%\marginnote{[Bonus]}
		\else
		\marginnote{\mbox{[\ExerciseTitle]\qquad}}
		\fi
		#2
	\end{Exercise}
	\columnbreak
		#3
	\end{multicols}
	\begin{Answer}[ref=#4]
		#5
	\end{Answer}
}





\SetDate[5/09/2026]

\begin{document}
\fancyhead[C]{\textbf{Équations — Exercices 1}}

\subsection*{Automatismes}


\exe{}{
	Trouver le $x\in\Q$ rationnel vérifiant chaque équation suivante, et l'exprimer sous forme de fraction irréductible.

	\begin{multicols}{2}
	\begin{enumerate}[itemsep=10pt]
		\item $3x + 5 = 2x + 7$
		\item $4x - 6 = 3x + 8$
		\item $5x + 10 = 2x + 4$
		\item $-2x + 3 = 4x - 1$
		\item $7x - 5 = 5x + 9$
		\item $6x + 4 = 3x + 12$
		\item $-3x + 2 = -5x - 4$
		\item $8x - 9 = 6x + 3$
		\item $2x + 11 = -3x + 5$
		\item $-4x + 7 = -x + 10$
	\end{enumerate}
	\end{multicols}
}{exe:inter-affine}
{
	\begin{multicols}{2}
	\begin{enumerate}[itemsep=10pt]
		\item $3x + 5 = 2x + 7 \iff x = 2$
		\item $4x - 6 = 3x + 8 \iff x = 14$
		\item $5x + 10 = 2x + 4 \iff x = -2$
		\item $-2x + 3 = 4x - 1 \iff x = \dfrac{2}{3}$
		\item $7x - 5 = 5x + 9 \iff x = 7$
		\item $6x + 4 = 3x + 12 \iff x = \dfrac{8}{3}$
		\item $-3x + 2 = -5x - 4 \iff x = -3$
		\item $8x - 9 = 6x + 3 \iff x = 6$
		\item $2x + 11 = -3x + 5 \iff x = \dfrac{-6}{5}$
		\item $-4x + 7 = -x + 10 \iff x = -1$
	\end{enumerate}
	\end{multicols}
}

\subsection*{Exercices}

\exe{, difficulty=1}{
	Pour chaque paire d'équations, ajouter entre elles le symbole
		\begin{enumerate}[itemindent=1.5cm]
			\item[$\implies$] si la première équation implique la seconde ;
			\item[$\impliedby$] si la deuxième équation implique la première ;
			\item[$\iff$] si les deux équations sont équivalentes ;
			\item[$\centernot\iff$] si les deux équations sont indépendantes.
		\end{enumerate}
	Spécifier l'opération appliquée pour obtenir l'équation impliquée à partir de l'autre.
	
	
	\begin{center}
	\begin{tabular}{ccc|c|c}
		Équation 1 & Symbole & Équation 2 & Opération & \thead{Les équations sont-\\ elles équivalentes ?} \\ \hline
		$3x + 2y = 1$ &  & $6x  + 4y = 2$ & &  \\ \hline
		$-x + 2y = 1$ &  & $-x + 2y - 1 = 0$ & &  \\ \hline
		$3x + 2y = 1$ & & $-3x - 5y = -3$ & &  \\ \hline
		$x^2 = -3x$ &  & $x = -3$ & &  \\ \hline
		$-10x - y = 3$ &  & $0=0$ & &  \\ \hline
		$x-y= 3$ & & $-x + y = -3$ & &  \\ \hline
		$x -y = 0$ & & $x=y$ & &  \\ \hline
		$-x + 2y = 3$ & & $0=0$ & &  \\ \hline
		$2x + 3y + 4 = 2x + 3y + 5$ & & $1=0$ & &  \\ \hline
		$-x - y = 1$ & & $-x - y = 1$ & &  \\ \hline
		$0=0$ & & $23x - 10y = 6$ & &  \\ \hline
		$x + 2y = 3$ & & $x^2 + 2yx = 3x$ & &  \\ \hline
		$x-3y  =0$ & & $-2x+6y = 0$ & &  \\ \hline
		$-2x + 3y = -3$ & & $4x - 6y = -6$ & &  \\ \hline
		$x^2 = x$ & & $x = 1$ & &  \\ \hline
		$x = 3$ & & $x^2 = 9$ & &  \\ \hline
	\end{tabular}
	\end{center}
	
	{Remarque : par abus de notation et pour gagner du temps, l'absence de symbole signifie en général que les équations sont équivalentes, mais il faut toujours faire attention à ce que cela soit bien le cas !}
	
	\textbf{Attention !} si $f(x) = 0 \implies g(x) = 0$, alors une solution de $f(x) = 0$ est une solution de $g(x) = 0$.
	En revanche, une solution de $g(x) = 0$ peut ne pas être une solution de $f(x) = 0$.
	
}{exe:1}{
	\begin{center}
	\begin{tabular}{ccc|c|c}
		Équation 1 & Symbole & Équation 2 & Opération & \thead{Les équations sont-\\ elles équivalentes ?} \\ \hline
		$3x + 2y = 1$ & {$\iff$} & $6x  + 4y = 2$ & {$\times2$} & {Oui} \\ \hline
		$-x + 2y = 1$ & {$\iff$} & $-x + 2y - 1 = 0$ & {$-1$} & {Oui} \\ \hline
		$3x + 2y = 1$ & {$\centernot\iff$} & $-3x - 5y = -3$ & & {Non} \\ \hline
		$x^2 = -3x$ & {$\impliedby$} & $x = -3$ & & {Non} \\ \hline
		$-10x - y = 3$ & {$\implies$} & $0=0$ & {$\times0$} & {Non} \\ \hline
		$x-y= 3$ & {$\iff$} & $-x + y = -3$ & {$\times(-1)$} & {Oui} \\ \hline
		$x -y = 0$ &{$\iff$} & $x=y$ & {$+y$} & {Oui} \\ \hline
		$-x + 2y = 3$ & {$\implies$} & $0=0$ & {$\times0$} & {Non} \\ \hline
		$2x + 3y + 4 = 2x + 3y + 5$ & {$\iff$} & $1=0$ & {$+(-2x-3y)$} & {Oui} \\ \hline
		$-x - y = 1$ & {$\centernot\iff$} & $-x - y = 1$ & & {Non} \\ \hline
		$0=0$ & {$\impliedby$} & $23x - 10y = 6$ & & {Non} \\ \hline
		$x + 2y = 3$ & {$\implies$} & $x^2 + 2yx = 3x$ & {$\times x$} & {Non} \\ \hline
		$x-3y  =0$ & {$\iff$} & $-2x+6y = 0$ & {$\times(-2)$} & {Oui} \\ \hline
		$-2x + 3y = -3$ & {$\centernot\iff$} & $4x - 6y = -6$ & & {Non} \\ \hline
		$x^2 = x$ & {$\impliedby$} & $x = 1$ & & {Non} \\ \hline
		$x = 3$ & {$\implies$} & $x^2 = 9$ & {Mise au carré} & {Non} \\ \hline
	\end{tabular}
	\end{center}
}

%\filbreak

\exe{difficulty=1}{
	Un élève de Seconde fournit la production écrite suivante.
	Celle-ci présente quelques erreurs.
	
	\fbox{\parbox{\textwidth}{
		\textbf{\oldunderline{Énoncé :}}
		
		Considérons l'équation $x^3 + x^2 + x + 1 = 0$ où $x\in\R$ est un nombre réel encore inconnu.
		\begin{enumerate}
			\item
			Montrer que $(x-1)(x^3 + x^2 + x + 1) = x^4 - 1$.
			\item
			En déduire la ou les solutions de l'équation originelle.
		\end{enumerate}
		
		\textbf{\underline{Réponse de l'élève :}}
		
		\begin{enumerate}
			\item
			Pour développer $(x-1)(x^3 + x^2 + x + 1)$,
			je distribue d'abord le $x$, puis le $-1$, et j'ajoute les résultats.
				\[ x(x^3 + x^2 + x + 1) + (-1)(x^3 + x^2 + x + 1) = x^4 + x^3 + x^2 + x -x^3 -x^2 - x - 1. \]
			Les puissances de $x$ égales se regroupent pour simplifier l'expression
				\[ x^4 + x^3 + x^2 + x -x^3 -x^2 - x - 1 = x^4 - 1, \]
			comme voulu.
			\item
			Un $x\in\R$ pour lequel $x^3 + x^2 + x + 1 = 0$ vérifie donc $0 = x^4 - 1$ et donc
			$x^4  = 1$.
			En appliquant la racine carrée, je trouve $x^2 = 1$ car $x^4 = (x^2)^2$.
			En l'appliquant encore, j'obtiens finalement $x = 1$.
			
			L'unique solution de l'équation est donc $1$.
		\end{enumerate}
	}}
	\begin{enumerate}
		\item 
		Identifier les erreurs commises par l'élève.
		\item
		Répondre à l'énoncé.
	\end{enumerate}
}{exe:erreurs}{
	\begin{enumerate}
		\item
		La première question est bien traitée et ne comporte aucune erreur.	
		\item
		La deuxième question contient deux erreurs : une lors de la résolution de $x^4 = 1$, et l'autre dans sa conclusion.
		
		Premièrement, appliquer la racine carrée est pertinente, mais $\sqrt{x^2} = |x|$, et non pas $x$ !
		La valeur absolue est importante pour ne pas oublier de solution.
		Ainsi, si $x^4 = 1$, alors $|x^2| = 1$. Comme $x^2 \geq 0$ pour tout $x\in\R$, la valeur absolue peut disparaître ici.
		Puis, si $x^2 = 1$, alors $|x| = 1$, et donc $x=1$ ou $x=-1$. Il y a donc deux solutions réelles à l'équation $x^4  = 1$.
		
		Deuxièmement, la conclusion est erronnée car elle présente une erreur de logique.
		En effet, l'impliquation $x^3 + x^2 + x + 1 = 0 \implies x \in \{-1 ; 1\}$ est vraie, mais l'implication inverse ne l'est pas !
		Nous disposons de deux candidats de solutions qu'il faut vérifier. La solution $x=1$ ne fonctionne pas, mais $x=-1$, elle, fonctionne.
		La seule solution est donc $x=-1$.
		
		En multipliant l'équation originelle par $x-1$, nous avons ajouté artificiellement la solution $x=1$.
		Par exemple, les implications suivantes sont vraie
			\[ x - 1 = 0 \implies (x-3)(x-1) = 0 \implies x \in \{3 ; 1\}. \]
	\end{enumerate}
}

\exe{}{
	Résoudre les équations suivantes pour $x\in\R$ non nul. 
	%Exprimer $x$ sous la forme $q \sqrt{n}$ avec $q\in\Q$ rationnel et $n\in\N$ le plus petit possible.
	\begin{multicols}{4}
	\begin{enumerate}
		\item $\dfrac{x}3 = \dfrac12.$
		\item $ \dfrac7x  =\dfrac12$
		\item $\dfrac6x  =\dfrac{\sqrt{3}}2$
		\item $\dfrac{x}2 =\dfrac1{\sqrt{3}}$
		%\item $\dfrac9x  =\dfrac{\sqrt{6}}5$
		%\item $-\dfrac3x =\sqrt{7}$
		%\item $-\dfrac7{2x}  =\dfrac{\sqrt{11}}3$
		%\item $\dfrac3{5x} =-\sqrt{5}$
	\end{enumerate}
	\end{multicols}
}{exe:aut0}{
	\begin{multicols}{2}
	\begin{enumerate}
		\item 
			\begin{align*}
				\dfrac{x}3 &= \dfrac12 \\
				x &= 3 \cdot \frac12 \\
				x &= \frac32
			\end{align*}
		\item
			\begin{align*}
				\dfrac7x  &=\dfrac12 \\
				7 &= x \cdot \frac12 \\
				x &= 7\cdot2 = 14
			\end{align*}
		\item
			\begin{align*}
				\dfrac6x  &= \dfrac{\sqrt{3}}2 \\
				6 &= x \cdot \dfrac{\sqrt{3}}2 \\
				x &= 6\cdot\dfrac2{\sqrt{3}} \\
				x &= \frac{12}3 \sqrt3 = 4\sqrt3
			\end{align*}
		\item
			\begin{align*}
				\dfrac{x}2 &= \dfrac1{\sqrt{3}} \\
				x &= 2 \cdot \dfrac1{\sqrt{3}} \\
				x &= \frac23 \sqrt3
			\end{align*}
	\end{enumerate}
	\end{multicols}
}

\exe{}{
	Trouver le $x \in \Q$ rationnel (et tel que chaque dénominateur est non nul) vérifiant chaque équation suivante, et l'exprimer sous forme de fraction irréductible.
	
	\begin{multicols}{3}
	\begin{enumerate}[itemsep=10pt]
		\item $9 = \dfrac3x$
		\item $-4 = \dfrac2x$
		\item $0,2422 = \dfrac1x$
		\item $0,7 = \dfrac{0,1}x$
		\item $3 - \dfrac4x = -9 + \dfrac1x$
		\item $-2 + \dfrac{-2}x = \dfrac3x - 10$
		\item $\dfrac1{x-2} = 3$
		\item $\dfrac3{x+5} = \dfrac7{x+5} - 3$
	\end{enumerate}
	\end{multicols}

}{exe:eq-lin}
{
	\begin{multicols}{2}
	\begin{enumerate}[itemsep=10pt]
		\item $9 = \dfrac3x \iff x = \dfrac13$
		\item $-4 = \dfrac2x \iff x=-\dfrac12$
		\item $0,2422 = \dfrac1x \iff x = \dfrac{10000}{2422} = \dfrac{5000}{1211}$
		\item $0,7 = \dfrac{0,1}x \iff x = \dfrac17$
		\item $3 - \dfrac4x = -9 + \dfrac1x \iff x = \dfrac5{12}$
		\item $\dfrac1{x-2} = 3 \iff x=\dfrac73$
		\item $\dfrac3{x+5} = \dfrac7{x+5} - 3 \iff x=\dfrac{-11}3$
	\end{enumerate}
	\end{multicols}
}


\subsection*{Approfondissements}

\exe{difficulty=1}{
	Montrer que les nombres suivants sont rationnels en les exprimant sous forme de fraction d'entiers.
	\[
	\begin{aligned}
		A &= 0,666{\dots} \\
		B &= 1,666{\dots} \\
	\end{aligned}
	\hspace{5cm}
	\begin{aligned}
		C &= 0,121212{\dots} \\
		D &= 0,34777{\dots} \\
	\end{aligned}
	\]
	\[
	E = 0,123456789123456789123{\dots} \text{ (nombre d'Amandine) }
	\]
	
	\emph{Indication : montrer d'abord que $10A = 6+A$.}
}{exe:dev-to-fraction}{
	\begin{enumerate}[label=\Alph*.]
		\item 
		La relation $10A = 6 + A$ donne $A = \frac23$.
		\item
		La relation $10B = 15 + B$ donne $B = \frac{15}{9} = \frac53$.
		On aurait aussi pû utiliser que $B = 1+A$.
		\item
		La relation $100C = 12 + C$ donne $C = \frac{12}{99} = \frac{4}{33}$.
		\item
		En posant $\tilde{E} = 100E = 34,777{\dots}$, on obtient $10\tilde{E} = 313 + \tilde{E}$, d'où $\tilde{E} = \frac{313}{9}$, et donc $E = \frac{1}{100}E' = \frac{313}{900}$.
		\item 
		La relation $10^9 E = 123~456~789 + E$ donne $E = \frac{123~456~789}{999~999~999}$.
	\end{enumerate}
}



%\exe{, difficulty=1}{
%	Multipliez l'équation 
%		\[ \dfrac{a}b + \dfrac{c}d = x \]
%	par $b$ puis $d$ à gauche et à droite et en déduire la règle d'addition des fractions :
%		\[ x = \dfrac{ad + cb}{bd}. \]
%}{exe:common-denom}{
%	En multipliant par $bd$, on trouve
%		\[ (bd) x =  bd \left( \dfrac{a}b + \dfrac{c}d \right) = ad + bc, \]
%	d'où
%		\[ x = \dfrac{ad+bc}{bd}. \]
%}


%
%\exe{}{
%	Résoudre les équations suivantes pour $x\in\R$ non nul. Exprimer $x$ sous la forme $q \sqrt{r}$ avec $q\in\Q$ rationnel et $r\in\N$ le plus petit possible.
%	\begin{multicols}{2}
%	\begin{enumerate}
%		\item $\dfrac3x = \dfrac73.$
%		\item $ \dfrac7x  =\dfrac1x + 2$
%		\item $\dfrac{16}x  =\dfrac{\sqrt{3}}2$
%		\item $-\dfrac2x =4\sqrt{3}$
%		\item $\dfrac9{-x}  =\dfrac{\sqrt{6}}5$
%		\item $-\dfrac3x =\dfrac57\sqrt{7}$
%		\item $-\dfrac9{2x}  =\dfrac{\sqrt{11}}3$
%		\item $\dfrac2{7x} =-\sqrt{14}$
%	\end{enumerate}
%	\end{multicols}
%}{exe:eq-sqrt}{}
%
%\exe{difficulty=1}{
%	\begin{multicols}{2}
%	La Terre complète son tour autour du Soleil en environ 365,2422 jours. 
%	Un calendrier ne peut donc pas contenir 365 jours tous les ans si celui-ci doit être synchronisé avec les saisons.
%	
%	Si le calendrier contient moins de 365,2422 jours, il sera en avance.
%	S'il contient plus de 365,2422 jours, il sera en retard
%
%	\begin{center}
%		\includegraphics[page=4]{../../CM/2-figures.pdf}
%	\end{center}
%	\end{multicols}
%	
%	\begin{enumerate}
%	\item
%	Calendrier julien\footnote{De Jules César (né en 100 av. J.-C., mort en 44 av. J.-C.).}.
%	
%	En $46$ av. J.-C., Jules César introduit la règle de l'année bissextile.
%	\begin{enumerate}
%		\item
%		Trouver un entier naturel $a \in \N$ non nul tel que
%			\[ 365{,}2422 \approx 365 + \dfrac1a.\]
%
%		\item
%		En déduire la règle d'année bissextile du calendrier julien.
%		
%		\item
%		Au bout de 1 600 ans, de combien de jours le calendrier est-il en retard ?
%	\end{enumerate}
%	
%	\item
%	Calendrier grégorien\footnote{Du pape Grégoire XIII (1502-1585).}.
%	
%	En 1582, plus de 1 600 ans après l'introduction de la règle d'année bissextile, le calendrier julien prend donc un retard notable.
%	Le pape Grégoire XIII promulgue alors une réforme du calendrier, ainsi créant le calendrier grégorien encore utilisé aujourd'hui.
%	
%	Afin de corriger le problème de dérive du calendrier, une nouvelle règle est ajoutée à la définition d'année bissextile.
%	
%	\begin{enumerate}
%		\item
%		Trouver un entier naturel $b \in \N$ tel que
%			\[ 365{,}2422 \approx 365 + \dfrac14 - \dfrac{b}{400}.\]
%
%		\item
%		En déduire la règle d'année bissextile du calendrier grégorien.
%		
%		\item
%		Au bout de combien d'années le calendrier est-il en retard d'un jour ?
%	\end{enumerate}
%	
%	\item
%	Un nouveau calendrier ?
%	\begin{enumerate}
%		\item
%		Trouver un entier naturel $c \in \N$ tel que
%			\[ 365{,}2422 \approx 365 + \dfrac14 - \dfrac{3}{400} - \dfrac{c}{4000}.\]
%
%		\item 
%		Créer une nouvelle règle d'année bissextile.
%		\item
%		Calculer le nombre d'années avant que votre calendrier soit en retard d'un jour.
%	\end{enumerate}
%	\end{enumerate}
%}{exe:calendrier}{}
%
%
%\exe{}{
%	On souhaite connaître les solutions de l'équation du deuxième degré d'inconnue $x\in\R$ :
%		\begin{align}
%			22x^2 - 125x + 22  = 0. \label{eq:1}
%		\end{align}
%	\begin{enumerate}
%		\item Montrer que $22x^2 - 125x + 22 = (2x-11)(11x-2)$.
%		\item En déduire l'ensemble des solutions de l'équation \eqref{eq:1}.
%	\end{enumerate}
%}{exe:eq7}{
%	\begin{enumerate}
%		\item On part toujours de la forme factorisée qu'on développe par double distributivité :
%			\[ (2x - 11)(11x - 2) = 22x^2 - 4x - 121x + 22 = 22x^2 - 125x + 22. \]
%		\item On utilise que le produit est nul si et seulement si un des deux facteurs est nul.
%		L'ensemble des solutions de \eqref{eq:1} est donc $\left\{ \dfrac{11}2 ; \dfrac2{11} \right\}$.
%	\end{enumerate}
%}
%
%
%\exe{}{
%	On souhaite connaître les solutions de l'équation du troisième degré d'inconnue $x\in\R$ :
%		\begin{align}
%			4x^3 - 6x^2 - 2x + 3  = 0. \label{eq:2}
%		\end{align}
%	\begin{enumerate}
%		\item Montrer que $4x^3 - 6x^2 - 2x + 3 = (2x-3) \left(2x^2-1\right)$.
%		\item En déduire l'ensemble des solutions de l'équation \eqref{eq:2}.
%	\end{enumerate}
%}{exe:eq8}{
%
%	\begin{enumerate}
%		\item On part toujours de la forme factorisée qu'on développe par double distributivité.
%		On utilise aussi que $x^2 \cdot x = x \cdot x \cdot x = x^3$.
%			\[ (2x - 3)(2x^2 - 1) = 4x^3 - 2x -6x^2 + 3. \]
%		\item On utilise que le produit est nul si et seulement si un des deux facteurs est nul.
%		Or l'équation $2x^2 - 1 = 0$ est équivalente à $x^2 = \dfrac12$.
%		En prenant la racine et en utilisant que $\sqrt{x^2} = |x|$, on obtient
%			\[ |x| = \sqrt{\dfrac12} = \dfrac1{\sqrt2}. \]
%		Les deux valeurs $\pm \dfrac1{\sqrt2}$ (notation $\pm$ pour \og plus ou moins \fg) sont donc solutions.
%		On en déduit que l'ensemble des solutions de \eqref{eq:2} est $\left\{ \dfrac32 ; \pm \dfrac1{\sqrt2} \right\} = \left\{ \dfrac32 ; \dfrac1{\sqrt2} ; - \dfrac1{\sqrt2} \right\}$
%	\end{enumerate}
%
%}
%
%
%\exe{}{
%	On souhaite connaître les solutions de l'équation du quatrième degré d'inconnue $x\in\R$ :
%		\begin{align}
%			16x^4 - 24x^2 + 9  = 0. \label{eq:3}
%		\end{align}
%	\begin{enumerate}
%		\item Montrer que $16x^4 - 24x^2 + 9 = \left(4x^2-3\right)^2$.
%		\item En déduire l'ensemble des solutions de l'équation \eqref{eq:3}.
%	\end{enumerate}
%}{exe:eq9}{
%
%	\begin{enumerate}
%		\item On part toujours de la forme factorisée qu'on développe à l'aide de l'identité remarquable $(a-b)^2 = a^2 + b^2 -2ab$.
%		On utilise aussi que $(x^2)^2 = x^2 \cdot x^2 = x\cdot x\cdot x\cdot x = x^4$.
%			\[ (4x^2 - 3)^2 = (4x^2)^2 + 3^2 - 2\cdot(4x^2)\cdot3 = 16x^4 + 9 -24x^2. \]
%		\item On utilise que le produit est nul si et seulement si un des deux facteurs est nul.
%		Ici, les deux facteurs sont les mêmes, car $(4x^2 - 3)^2 = (4x^2 - 3)(4x^2 - 3)$, donc on se remet à résoudre $(4x^2 - 3) = 0$.
%		Ceci est équivalent à $x^2 = \dfrac34$, et donc l'ensemble des $x$ vérifiant \eqref{eq:3} est donné par $\left\{ \pm \sqrt{\dfrac34} \right\} = \left\{ \pm \dfrac12\sqrt{3} \right\}= \left\{ \dfrac12\sqrt{3} ;  -\dfrac12\sqrt{3} \right\}$.
%	\end{enumerate}
%
%}
%
%
%
%
%
%\exe{}{
%	Montrer que les nombres suivants sont rationnels en les exprimant sous forme de fraction d'entiers.
%	\begin{align*}
%		A = 0,123123123{\dots} &&
%		B = 0,9999{\dots}
%	\end{align*}
%}{exe:dev-to-fractionE}{
%	\begin{enumerate}[label=\Alph*.]
%		\item
%		La relation $1000A = 123 + A$ donne $A = \frac{123}{999} = \frac{41}{333}$.
%		\item 
%		La relation $10B = 9 + B$ donne $9B = 9$ et $B=1$.
%	\end{enumerate}
%}
%
%
%\exe{, difficulty=2}{
%	Donner quelques termes en plus de la somme infinie suivante. 
%		\[ \dfrac1{10} + \dfrac1{100} + \dfrac1{1 000} + \dfrac1{10 000} + \dots \]
%	Écrire le nombre en écriture décimale, puis l'exprimer sous forme de fraction de deux entiers comme à l'exercice \ref{exe:dev-to-fraction}.
%}{exe:20}{
%	La somme infinie vaut $0,1111{\dots} = \frac19$.
%}
%
%\exe{, difficulty=2}{
%	Donner quelques termes en plus de la somme infinie suivante.
%		\[ \dfrac12 + \dfrac14 + \dfrac18 + \dfrac1{16} + \dots \]
%	À quoi est-elle égale ? Un résultat entier est attendu.
%}{exe:21}{
%	Posons $S$ cette somme. 
%	En multipliant par $2$, on trouve $2S = 1 + S$, et donc $S=1$.
%}
%
%\exe{, difficulty=2}{
%	On considère la somme finie suivante
%		\[ S = 1 + 2 + 4 + 8 + 16 + \dots + 2^{n-1} +  2^n, \]
%	où $n\in\N$ est un entier naturel.
%	
%	Montrer que $2S = S - 1 + 2^{n+1}$, et en déduire la valeur de $S$ sous forme close (sans points de suspension).
%}{exe:22}{
%	En multipliant par 2, le premier terme (1) disparait, et un nouveau terme ($2^{n+1}$) apparaît.
%	On a donc bien $2S = S -1 + 2^{n+1}$, et donc $S = 2^{n-1} - 1$.
%}
%
%
%\exe{, difficulty=1}{
%	Montrer qu'un triangle isocèle rectangle ne peut pas avoir des côtés de longueurs toutes entières.
%}{exe:non-iso-rec}{
%	Supposons qu'il existe un triangle isocèle rectangle de côtés $(a ; a; b)$.
%	
%	Or, comme le triangle est rectangle isocèle,
%		\[ b^2 = a^2 + a^2 = 2a^2 \iff b = \pm\sqrt2 a. \]
%	Par conséquent, $\sqrt2 = \pm\frac{b}{a} \in \Q$, ce qui contredit le théroème d'irrationalité de la racine carrée de 2.
%}
%
%\exe{, difficulty=1}{
%	Soit $q\in\Q$ un rationnel. Montrer qu'un triangle dont les côtés sont de longueur $(4 ; 5 ; q)$ ne peut être rectangle que si $q=3$.
%	
%	En autorisant $x\in\R$ irrationel, donner un triangle rectangle de côtés $(4 ; 5; x)$ avec $x\neq3$.
%}{exe:non-rec}{
%	Il y a deux cas de figure à étudier : $q$ est la longueur de l'hypothénuse, ou elle ne l'est pas.
%	Dans le premier cas, le théorème de Pythagore énonce que
%		\[ 4^2 + 5^2 = q^2 \iff 41 = q^2. \]
%	Il suit que $q=\sqrt{41}$ qui n'est pas rationnel car 41 n'est pas un carré parfait.
%	Ceci répond à la deuxième question : $(4 ; 5; \sqrt{41})$ est rectangle.
%	
%	Dans le second cas,
%		\[ q^2 + 4^2 = 5^2 \iff q^2 = 9. \]
%	Par conséquent, $q = \pm3$ et seule la solution $q=3$ est positive et donc une longueur valide.
%}

%%%%%%%%%%%

\newpage
\fancyhead[C]{\textbf{Nombres rationnels — Solutions}}
\shipoutAnswer

\end{document}
